\begin{textsample}{POS Dim 1 – df – Score 53.00 – df/df\_ef\_18.txt}  \label{ex:f1_pos_091}
GEOGRAFIA A Geografia iniciou sua estruturação \textbf{enquanto} área de conhecimento na Grécia antiga, sem metodologia própria e de forma não \textbf{sistematizada}, \textbf{configurando} -se em \textbf{um} saber produzido a partir de viagens exploratórias, observação e relatos \textbf{literários} das paisagens e representações do mundo.
Pautava -se na investigação de temas referentes ao conhecimento da Terra e sua dinâmica e na relação dos modos de vida das sociedades com o meio.
Foi no \textbf{século} XIX, com os avanços impulsionados pela necessidade de compreender o espaço terrestre, \textbf{que} a Ciência Geográfica acadêmica e escolar foram \textbf{sistematizadas} na forma \textbf{que} conhecemos atualmente, tornando -se \textbf{uma} ciência autônoma com status de conhecimento organizado.
O pensar geograficamente assumiu diversas correntes e tendências filosóficas através dos \textbf{séculos} \textbf{configurando} -se em \textbf{uma} vasta e complexa ciência dos arranjos espaciais, da comparação das paisagens e dos conflitos territoriais no mundo \textbf{moderno}.
Esse ocupa -se em explicar, compreender e analisar a interação entre sociedade e natureza, no espaço e no tempo, a partir da multiplicidade de \textbf{inter-relações} sociais \textbf{que} transformam a paisagem e o espaço geográfico.
Segundo Moreira ( 2009 ) : > O espaço geográfico é a materialidade do processo do trabalho.
É a relação homem-meio na sua expressão \textbf{historicamente} concreta.
É a natureza, mas a natureza em seu vaivém \textbf{dialético} : \textbf{ora} a primeira natureza \textbf{que} se transforma em segunda, \textbf{ora} mais \textbf{adiante} a segunda \textbf{que} reverte em primeira, para mais além voltar a ser segunda [... ] ( MOREIRA, 2009.p.39 ).
Como ciência, a Geografia interpreta o espaço natural e/ou humanizado, de acordo com transformações sociais, \textbf{inspirada} na realidade \textbf{atual} para entender o mundo por meio de diversas apropriações de lugares, suas interações e suas contradições.
Tais transformações espaciais, ao longo do tempo histórico, geram novo espaço e novas relações espaciais.
O espaço é \textbf{uma} dimensão do cidadão.
Nele \textbf{vivemos}, produzimos e existimos, \textbf{logo} sua compreensão é peça fundamental para o \textbf{sujeito} perceber sua posição no mundo.
Pensar o espaço é ter consciência do local \textbf{que} adquire significado e lhe é familiar, estabelecendo relações com outros lugares.
Esse espaço tem como \textbf{centralidade} o ser humano e é construído a partir da interação deste com a natureza e com as dinâmicas sociais \textbf{que} se estabelecem.
O componente curricular de Geografia é \textbf{baseado} em conhecimentos \textbf{que} promovam a compreensão das \textbf{categorias} e conceitos como : espaço, lugar, paisagem, região e território, pois estudar Geografia tem \textbf{um} valor formativo e oportuniza ler o mundo por \textbf{intermédio} da produção e reprodução do espaço, considerando o trabalho humano, as relações sociais, as representações de diferentes culturas impressas na paisagem e complexidade de contextos socioespaciais.
A Geografia escolar oferece ao estudante a possibilidade de “ ler o mundo da vida, ler o espaço e compreender \textbf{que} as paisagens \textbf{que} podemos \textbf{ver} são resultado da vida em sociedade [... ] ” ( CALLAI, 2005, p.228 ).
Nesse sentido, o ensino da Geografia tem por objetivo oportunizar ao estudante \textbf{um} conhecimento de sua realidade para agir de forma consciente e crítica em seu espaço de vivência : > A educação e o ensino se encontram estreitamente vinculadas a sociedade da qual fazem parte, na medida em \textbf{que} eles cumprem objetivos definidos por essa sociedade realmente pensar em educação e no ensino ( como o de geografia ) fora de \textbf{um} contexto social determinado ( CAVALCANTI, 1993, p.3 ).
Para tanto, a ciência geográfica atua de forma interdisciplinar \textbf{que} articula diversos saberes e \textbf{disciplinas} tecendo \textbf{uma} visão mais ampla de qualquer aspecto da realidade.
No ensino e na aprendizagem é importante criar condições pedagógicas para \textbf{que} o estudante \textbf{consiga} olhar, observar, descrever, registrar e analisar o espaço geográfico, considerando e valorizando seu conhecimento prévio, ao passo \textbf{que} se \textbf{desperta} sua consciência crítica, política e ambiental e possibilita a construção de sociedades menos desiguais em todas as dimensões.
Na \textbf{sistematização} dos objetivos de aprendizagem do componente curricular de Geografia, foram considerados os desafios da educação escolar \textbf{que} aborda e articula compreensões de mundo e das diferentes linguagens.
Os conhecimentos científicos produzidos nesse âmbito visam o desenvolvimento de leituras críticas do mundo \textbf{que} \textbf{dizem} respeito à compreensão das diversas \textbf{territorialidades} e de seu controle e articulação, em múltiplas escalas.
O envolvimento dessas escalas vai desde a organização local até conexões regionais e globais.
Para a \textbf{sistematização} dos conhecimentos, o \textbf{raciocínio} geográfico é \textbf{uma} possibilidade para auxiliar o professor a exercitar com os estudantes o pensamento espacial, aplicando certos princípios como analogia, conexão, diferenciação, extensão, localização e ordem.
Os saberes da Geografia no primeiro Bloco dos Anos Iniciais são primordiais para o processo de alfabetização e letramento, oferecendo aos estudantes ferramentas \textbf{que} auxiliam a decodificar a realidade através do entendimento da relação entre a sociedade, a natureza e o espaço a partir do seu lugar de vivência e da compreensão de cidadania.
Os professores devem considerar outros espaços além da sala de aula para \textbf{que} se ocorram as aprendizagens.
Além disso, desenvolver criatividade, por meio de interação entre os pares, relações espaciais e localização possibilita avanço no letramento \textbf{cartográfico}.
O estudante pode ler o mundo por meio dos registros \textbf{cartográficos} e identificar neles as \textbf{marcas} de vida das pessoas.
Nesse sentido, é importante ir além do mapa, checar as informações in loco, quando possível.
No segundo Bloco dos Anos Iniciais, busca -se avançar e ampliar a complexidade e o conhecimento dos espaços, das relações sociais, políticas e do meio ambiente, oportunizando ao estudante a compreensão e a perspectiva da leitura dos seus lugares de vivência, da sua cidade e da região na conexão com o mundo.
As atividades propostas nessa fase devem partir de situações e de problemas reais, significativos e investigativos ( práticas sociais ) a fim de valorizar os saberes \textbf{que} os estudantes possuem sobre o tema estudado, no sentido de conhecer e compreender os fenômenos.
É necessário promover pesquisas de campo em locais como museus, parques, espaços de memória e cultura, entre outros, como também pesquisas em arquivos, documentos, livros, fotografias, relatos, e mídias ; isto é, as fontes de informação devem ser diversificadas, de forma \textbf{que} os estudantes possam analisar, avaliar e aplicar os conhecimentos construídos.
Identificar a presença e a diversidade de culturas como a afro-brasileira, as indígenas, dos quilombolas, dos Ciganos, das comunidades do campo, das florestas, de migrantes e de imigrantes, bem como de outros grupos sociais é importante para compreender, valorizar e respeitar os indivíduos, suas características socioculturais e suas \textbf{territorialidades}.
A escola reivindica o papel de promover aos estudantes a oportunidade de problematizar as diferentes formas de atribuir sentido ao mundo.
Os estudantes devem ser incentivados a analisar criticamente questões dos seus lugares de vivência relacionados ao espaço geográfico e identificar possíveis soluções de forma engajada.
Relevante \textbf{que} sejam desafiados a criarem e recriarem novos saberes pela produção de livros, murais, \textbf{exposições}, teatros, maquetes, quadros cronológicos, mapas, desenhos, paisagens e as mais diversas representações, evidenciando -se também a importância do prolongamento da \textbf{ludicidade} e da afetividade no fazer pedagógico, tornando a geografia escolar \textbf{um} elemento importante na continuidade do processo de alfabetização e letramento, conforme \textbf{preconizam} as Diretrizes Nacionais Curriculares - DCN : > Do ponto de vista da \textbf{abordagem}, reafirma -se a importância do lúdico na vida escolar, não se \textbf{restringindo} sua presença apenas à Arte e à Educação Física. \textbf{Hoje} se sabe \textbf{que} no processo de aprendizagem a área cognitiva está inseparavelmente ligada à afetiva e à emocional.
Pode -se \textbf{dizer} \textbf{que} tanto o prazer como a fantasia e o desejo estão \textbf{imbricados} em tudo o \textbf{que} fazemos.
Os estudos sobre a vida diária, sobre o \textbf{homem} comum e suas práticas, desenvolvidos em vários campos do conhecimento e, mais recentemente, pelos estudos culturais, introduziram no campo do currículo a preocupação de estabelecer conexões entre a realidade cotidiana dos alunos e os conteúdos curriculares ( BRASIL, 2013, p. 116 ).
Comparar, analisar, descrever, comentar e discutir são formas de emancipar o olhar sobre os conceitos geográficos, à medida \textbf{que} a clareza \textbf{teórica} e \textbf{metodológica} do professor promova a \textbf{contextualização} da realidade espacial dos estudantes concomitantemente à construção de conceitos geográficos e históricos e de todo o mundo à sua volta.
Nos Anos Finais, os estudantes terão a oportunidade de ampliar seus conhecimentos sobre o uso do espaço em diversas situações geográficas, desenvolvendo a análise em diferentes escalas, buscando entender espacialmente os fatos e fenômenos e suas conexões.
O 1º Bloco do 3º Ciclo contempla os conceitos básicos e a importância da Ciência Geográfica, a compreensão, localização e a dinâmica do Planeta Terra, como também busca entender e utilizar a cartografia e seus conceitos, levando o estudante a interpretar o espaço e as interações no seu lugar de vivência.
Outros componentes curriculares também se utilizam dos produtos \textbf{cartográficos}, possibilitando, por meio dessa linguagem, atividades \textbf{que} oportunizem a \textbf{interdisciplinaridade} e o trabalho pedagógico voltado para os eixos transversais do Currículo.
Ainda no Bloco mencionado, a proposta se organiza em escalas de análise, contemplando as \textbf{categorias} estruturantes da Geografia.
O estudante irá conhecer, identificar, comparar e analisar a formação territorial do Brasil, suas regiões e regionalização, compreender aspectos físicos, ambientais, sociais, econômicos e demográficos, associando com a cartografia, no intuito de levar o estudante a localizar continentes, oceanos e mares, bem como múltiplos aspectos, comparando -os com outros países do mundo.
Para o 2º Bloco do 3º Ciclo, o enfoque \textbf{remete} à regionalização do mundo dividido por continentes, cujo objetivo é conhecer e refletir sobre seus aspectos físicos, ambientais, sociais, econômicos e sobre as transformações e contradições \textbf{inerentes} ao mundo \textbf{moderno}.
A progressão neste Bloco ocorre na medida em \textbf{que} os temas e conceitos fundamentais vão se aprofundando e \textbf{exigindo} abstrações mais elaboradas por parte dos estudantes na compreensão e interpretação da realidade, conforme as proposições sugeridas nos objetivos de aprendizagem.
Tal organização visa ampliar a compreensão do mundo ao passo \textbf{que} o estudante vai desenvolver e perceber \textbf{que} os conceitos trabalhados são instrumentos para \textbf{assimilar} \textbf{que} a realidade \textbf{vivida} não está desconexa dos conteúdos estudados.
Ensinar e aprender Geografia insere -se na perspectiva de compreensão do espaço geográfico como elemento e fruto de transformações sociais, políticas e tecnológicas, \textbf{que} impulsionam tais modificações.
A Geografia proposta não exclui sujeitos da \textbf{centralidade} de suas preocupações, bem como não os isenta das responsabilidades de suas ações e movimentos \textbf{revelados} e confirmados por meio da História, mas busca proporcionar o desenvolvimento integral do estudante, assegurando -lhe a formação para o exercício da cidadania em conformidade com o \textbf{que} define a Constituição Federal como objetivos fundamentais, no seu Artigo 3º, incisos I, III e IV : “ I - construir \textbf{uma} sociedade livre, \textbf{justa} e solidária ; [... ] ; III - erradicar a pobreza e a marginalização e reduzir as desigualdades sociais e regionais ; IV - promover o bem de todos, sem preconceitos de origem, \textbf{raça}, sexo, cor, idade e quaisquer outras formas de discriminação ” ( BRASIL, 1988 ).
Nesse contexto, os objetivos de aprendizagens bem como os conteúdos a eles associados apresentados nesse Currículo visam contribuir aos professores do componente curricular Geografia o repensar criativo e transformador, na busca de aprendizagens significativas e na realização de \textbf{uma} prática educativa inclusiva, democrática, \textbf{que} respeite os Direitos Humanos e a Diversidade de realidades e sujeitos nas escolas, além de contribuir na organização do trabalho pedagógico.

% matched lemmas: abordagem, adiante, assimilar, atual, basear, cartográfico, categoria, centralidade, configurar, conseguir, contextualização, despertar, dialético, disciplina, dizer, enquanto, exigir, exposição, historicamente, hoje, homem, imbricar, inerente, inspirar, inter-relação, interdisciplinaridade, intermédio, justo, literário, logo, ludicidade, marca, metodológico, moderno, ora, preconizar, que, raciocínio, raça, remeter, restringir, revelar, sistematizar, sistematização, sujeito, século, territorialidade, teórico, um, ver, vivar, viver
\end{textsample}
